\paragraph{Summary of empirical findings}
The calibration on 295{,}637 FEMA claims and the two application studies yield a coherent set of findings. The duration-dependent Richards function captures a 2.6 times damage amplification at 1~m depth between instantaneous and 288-hour flooding, operating through a single mechanism (inflection point shift, $\lambda_a = 0.49$). This amplification propagates through the loss pipeline: for ADM, the duration effect increases the total CapEx loss estimate by 3--13 percentage points per facility cluster compared to JRC functions. The claims-calibrated sector adjustment $\kappa$ further modulates the signal, preventing false alarms for diversified companies while preserving sensitivity for concentrated exposures.

Flood depth remains the dominant predictor of damage at the individual-claim level (74\% of gradient boosting feature importance; Table~\ref{tab: feature enrichment}), and the duration extension adds predictive value primarily at the aggregated and portfolio scale. At the claim level, both the depth-only Richards model and richer ML models leave an irreducible noise floor of $\approx 0.22$ MAE, reflecting unobservable idiosyncratic factors. The practical implication is that the duration-dependent formulation is best suited to event-level screening and portfolio risk ranking --- the use case this paper targets --- rather than individual asset loss estimation.

\paragraph{Claims-based $\kappa$ calibration}
The two-stage $\kappa$ calibration -- building-level vulnerability from 82{,}650 NFIP non-residential claims, combined with a coverage-gradient extrapolation to corporate scale -- addresses the principal limitation of prior approaches that relied on small company samples. The resulting $\kappa_{\text{total}}$ values (0.14--0.21 depending on building stories) are validated out-of-sample on 24 corporate flood events with median ratio $0.90\times$ (mean $|\!\log(\text{ratio})| = 0.37$). The storey-based discrimination is empirically confirmed: two-storey companies show tighter prediction accuracy than single-storey companies, consistent with the physical mechanism (fraction of building above water).

\paragraph{Indirect loss quantification}
Frameworks that capture only direct (CapEx) losses underestimate total credit impact when operational disruption is material. In the ADM analysis, the implied OpEx-to-CapEx ratio of 0.18 is consistent with World Bank post-disaster recovery studies~\cite{hallegatte2015indirect} and confirms the importance of temporal loss structure for accurate credit risk assessment. For the 24-company validation, indirect losses are implicitly embedded in reported disclosures, making direct decomposition infeasible, but the magnitude of $\kappa_{\text{total}}$ partly reflects these indirect effects.

\paragraph{Comparison with CLIMADA}
When calibrated to the same damage assumptions, our expected loss estimates are consistent with CLIMADA-style calculations (within $\sim$5\% deviation), confirming the validity of our damage assessment module. The key advantage remains that our formulation introduces inundation duration directly into the vulnerability function. This allows prolonged events such as Harvey-like floods to be distinguished from shallower or faster-receding events without changing the surrounding hazard framework.

\paragraph{Limitations}
The coverage-gradient factor $\rho_{\text{gradient}} = 0.25$ is extrapolated from the log-linear trend within NFIP coverage bands; this extrapolation to \$5\,M has not been independently validated. The FEMA calibration is based on US residential and commercial flood claims; transferability to different construction standards, flood mechanisms, and geographies requires validation. The OpEx model parameters are not independently calibrated but are grounded in post-disaster recovery studies.

The duration parameter $\tau$ in the current calibration is an event-level proxy, not a property-level observation. Event-level $\tau$ aggregates over substantial intra-event heterogeneity: during Harvey, for instance, FEMA flood extent reports indicate that some areas experienced inundation lasting fewer than 72~hours while others exceeded 14~days. Our three calibration events span an 8 times range in event-level $\tau$ (36--288~hours), which supports identification of the inflection-shift parameter $\lambda_a$, but the model is not validated at sub-event spatial resolution. Future work linking property-level claims to gauge-based flood duration maps would allow a direct test of whether intra-event $\tau$ variation drives additional explanatory power beyond depth alone.

This paper also does not model flow velocity as a separate intensity measure. Velocity contributes to structural damage through hydrodynamic loading (drag $\propto v^2$ and momentum flux $h \cdot v^2$), especially in riverine and coastal settings. Engineering studies of flash-flood and torrential events report that velocity-driven loads can add 30--80\% to depth-only damage estimates at $v \gtrsim 2$\,m\,s$^{-1}$~\cite{santangelo2021landslides, fahad2020coupled, fahad2022decision}; for the slower, predominantly stagnant flooding that characterises Sandy, Harvey, and Katrina, typical velocities are below 1\,m\,s$^{-1}$ and the velocity contribution is materially smaller. We therefore expect the depth-only Richards formulation to underestimate damage primarily in fast-moving riverine and dam-break floods; for such events an extension incorporating velocity (as in~\cite{merz2010review}) or a multivariate intensity measure is recommended.

\paragraph{Identification of $\tau$ vs.\ confounded covariates}
A related concern is that event-level $\tau$ may co-vary with unobserved covariates such as the depth distribution, flow velocity, antecedent moisture, and building-stock composition across the three hurricanes. The empirical $\lambda_a = 0.49$ should therefore be read as the partial effect of $\tau$ \emph{conditional on depth bins} but not strictly orthogonal to event-specific construction-quality and exposure-mix differences. We mitigate this concern in three ways: (i) calibration is performed on binned depth-damage data (12 bins per event) so that within-bin depth variation is absorbed; (ii) the $w_{\text{USACE}} = 3$ instantaneous anchor pins the $\tau = 0$ end of the curve to a duration-independent reference; and (iii) the leave-one-event-out test (Table~\ref{tab: loo results}) shows that the duration coefficient extrapolates to a held-out event, which would not occur if $\lambda_a$ were absorbing only event-specific noise. A fully causal identification would require pairing claims with hydraulic-model output containing depth, velocity, and local recession time at the property level; this is the natural next step but is not feasible at the current NFIP scale.

\paragraph{Predictive use and how to assign $\tau$ in practice}
The duration parameter is applied to future flood scenarios by reading $\tau$ from one of three sources, in decreasing order of preference: (a)~gauge-based flood-recession curves (e.g.\ NOAA AHPS or USGS streamflow records) for the nearest hydrological reference point, integrated above the floor elevation of the asset; (b)~hydrodynamic simulation output (e.g.\ HEC-RAS, LISFLOOD) where available; or (c)~a literature/event-report estimate of typical inundation duration for the relevant flood archetype (storm surge, riverine, pluvial). The framework therefore does not require that future events match the calibration events -- it only requires that an event-scale $\tau$ value can be assigned by the practitioner. We recommend treating $\tau$ as bounded between 24\,h (rapid drainage) and 720\,h (extreme stagnant flooding) and reporting damage estimates at $\tau$ and $\tau \pm 50\%$ as a sensitivity range when the duration estimate is uncertain (Section~\ref{ssec: tau sensitivity} for the ADM example).

\paragraph{Parameter uncertainty}
The Richards parameters in Table~\ref{tab: calibrated params} are point estimates from differential-evolution optimisation. To probe within-sample sensitivity, we re-ran the calibration on bootstrap replicates of the binned data (perturbing each bin mean by a Gaussian draw with the bin standard error and re-optimising). Across 200 replicates the qualitative ranking is preserved: $\lambda_a$ is positive and bounded away from zero in every replicate, while $\lambda_Q$ and $\delta_\nu$ remain close to zero. The objective is multi-modal in the $(\lambda_Q, \delta_\nu)$ corner; the bootstrap therefore reads as a local sensitivity analysis around the preferred mode rather than a global posterior. The non-significance of $\lambda_Q$ and $\delta_\nu$ should be interpreted as ``cannot be distinguished from zero given the optimiser bounds and three event-level observations'', not as a strict point identification of zero. The dispersion of bootstrap parameter estimates is small relative to the cross-event MAE differences in Table~\ref{tab: loo results}. We none the less emphasise that bootstrap dispersion does not quantify the dominant source of uncertainty, which is having only three distinct event-level $\tau$ values; reducing this requires expanding the calibration set rather than resampling within it. Full bootstrap replicates and binning-sensitivity outputs are available in the reproducibility code.

Although bootstrap confidence intervals are reported, the calibration rests on binned event summaries with only three non-zero duration levels; a fuller hierarchical uncertainty treatment remains a logical next step.
